\begin{center}
	\chapter{LITERATURE REVIEW}
\end{center}
\justifying

\section{Background and Related Works}

\subsection{Linear recurrent architectures}

Attention computes, at every position, a weighted combination over the entire
preceding sequence, giving cost quadratic in sequence length \cite{attention}.
Linear attention replaces the softmax with a kernel feature map, which allows the
summation to be reorganised into a recurrence over a fixed-size state and reduces
the cost to linear \cite{lineartransformers}. Schlag et al. \cite{fastweight}
subsequently identified this recurrence as a form of fast weight programming and
observed that the elementary additive update overwrites capacity as the sequence
lengthens, motivating an update rule with an explicit removal term.

Structured state-space models arrive at a similar recurrence from a different
direction, parameterising a linear dynamical system whose transition operator is
diagonal or diagonal-plus-low-rank \cite{s4}, with selective, input-dependent
parameterisations following later \cite{mamba}. Gated linear attention
\cite{gla} introduces data-dependent decay into the linear-attention recurrence
and supplies hardware-efficient chunked kernels for training it at scale.

\subsection{State tracking and its limits}

The capability that separates these architectures is state tracking. Merrill et
al. \cite{illusionofstate} show that state-space models whose transition operator
is diagonal cannot express the composition of non-commuting operations, placing
them below the complexity class required for permutation tracking regardless of
depth or width. Shakerinava et al. \cite{diagonalssm} refine this
characterisation for input-dependent complex-valued diagonal models, and report
in addition that multi-layer models frequently fail to learn state tracking for
non-Abelian groups even in configurations where they possess the representational
capacity to express it, indicating a gap between expressivity and optimisation.
Liu et al. \cite{shortcuts} establish the complementary result for
attention-based models, showing that shallow transformers learn shortcut
solutions to automata that generalise poorly beyond the training length.

Permutation-group word problems are the standard instrument in this line of work.
Sequences are drawn from an alphabet of permutations and the target at each
position is the running product. The symmetric group $S_5$ is the usual hard
case, since it is non-solvable and therefore lies outside what a constant number
of diagonal layers can express.

\subsection{Householder products as transition operators}

DeltaNet \cite{deltanet} makes the delta-rule recurrence parallelisable over
sequence length by expressing a chunk of updates as a matrix product, which
yields a state-transition operator that is a product of rank-one modifications of
the identity. DeltaProduct \cite{deltaproduct} generalises this by taking $n_h$
such factors per token, each a generalized Householder transformation
$H(\beta,k) = I - \beta k k^{\!\top}$ with $\lVert k \rVert = 1$ and
$\beta \in [0,2]$. The parameter $\beta$ interpolates between three regimes: at
$\beta = 0$ the factor is the identity, at $\beta = 1$ it is the orthogonal
projection removing the component along $k$, and at $\beta = 2$ it is the
reflection in the hyperplane orthogonal to $k$. Since a transposition of two
coordinates is precisely such a reflection, a permutation expressible as a
product of $\ell$ transpositions is realisable by $\ell$ reflection factors, and
DeltaProduct accordingly reports that state-tracking performance on $S_n$
improves as $n_h$ is raised towards the maximum transposition length. The number
of factors is a global hyperparameter, and the paper identifies per-token
adaptive order as an open direction.

\subsection{Adaptive computation}

Allocating a variable amount of computation per input is a well-established
mechanism. Adaptive Computation Time \cite{act} attaches a halting unit to a
recurrent step, accumulates halting probabilities until they exceed a threshold,
and penalises the expected number of steps through a ponder cost. PonderNet
\cite{pondernet} reformulates the halting distribution probabilistically and
supplies an unbiased gradient estimator with improved stability. Universal
Transformers \cite{universaltransformer} apply the same principle to depth in an
attention model. These mechanisms are mature and are adopted here without
modification; the contribution of the present work lies elsewhere.

\section{Outcome of Literature Review}

Table \ref{table:survey} summarises the surveyed works along the two axes
relevant to this project: whether the architecture can express non-commutative
state transitions, and whether the computation it spends per token is fixed or
variable.

\begin{table}[ht]
	\caption{Position of the surveyed architectures} % title of Table
	\centering
	\begin{tabular}{P{4.6cm} P{3.4cm} P{3.4cm}}
		\hline\hline
		Architecture family & Non-commutative state tracking & Computation per token \\ [0.5ex]
		\hline
		Softmax attention \cite{attention} & yes, at quadratic cost & fixed \\
		Linear attention \cite{lineartransformers, gla} & no & fixed \\
		Diagonal state-space models \cite{s4, mamba} & no \cite{illusionofstate} & fixed \\
		Delta rule \cite{fastweight, deltanet} & limited, one factor & fixed \\
		Householder products \cite{deltaproduct} & yes, for $n_h$ sufficiently large & fixed, global $n_h$ \\
		Adaptive computation \cite{act, pondernet} & not applicable & variable \\ [1ex]
		\hline
	\end{tabular}
	\label{table:survey}
\end{table}

Three observations follow from the survey.

First, the expressivity question for this architecture family is settled in
outline. Diagonal transition operators cannot track non-commutative structure,
products of rank-one modifications can, and the factor count governs how much
structure is reachable. What has not been established is the correct value of
that count for a given task, as distinct from the observation that larger values
perform better.

Second, the cost accounting in this line of work treats the required expressivity
as an attribute of the individual symbol, through the identification of the
factor count with the transposition length of the corresponding permutation. This
identification is exact for the standard representation of a symmetric group, but
it is not exact in general, and the discrepancy is not incidental. Chapter 3
establishes that the requirement is determined by the group generated by the
whole alphabet, and that the same symbol may require two factors under one
alphabet and four under another.

Third, adaptive computation supplies a mature halting mechanism that can be
applied without modification. The open question is therefore not how a model
should halt, but what the correct allocation is and whether it can be inferred
from the task without supervision.

\section{Problem Statement}

In linear recurrent architectures whose state-transition operator is a product of
generalized Householder factors, the number of factors per token, $n_h$, is a
global hyperparameter applied uniformly to every input symbol. The expressivity
actually required varies across symbols, so a model with a fixed order must be
provisioned for the maximum requirement and incurs that cost at every position.
Furthermore, the required value is not determined by local properties of the
symbols and cannot be obtained by inspection of the data. The problem addressed
in this work is to determine the correct per-token requirement, and to make the
factor count per-token and learned, without loss of accuracy relative to a
fixed-order model provisioned for the worst case.

\section{Objectives of the Project}

\begin{enumerate}[label={(\arabic*)}]
	\item To establish the correct expressivity cost model for state transitions
	expressed as products of generalized Householder factors, and to verify it by
	explicit construction rather than by optimisation.
	\item To reproduce the published state-tracking behaviour of the base
	architecture on the available hardware, so that all subsequent comparisons
	rest on a validated baseline.
	\item To design and implement a monotone halting gate that predicts a
	per-token factor count and satisfies the reachability constraints imposed by
	the expressivity analysis.
	\item To train the gate end to end under a penalty on expected computation,
	with no supervision on the allocation, and to compare the allocation it
	produces against the allocation derived from ground truth.
	\item To quantify the reduction in required computation that a correct
	per-token allocation achieves relative to the cheapest fixed order that solves
	the same task.
	\item To characterise the training behaviour of the architecture across
	alphabets of differing requirement distributions, and to establish a
	measurement protocol appropriate to that behaviour.
\end{enumerate}
\null
